\subsection{Predicates and Destination Images}

For an element size of $E$ bytes, logical lane $i$ is governed by predicate
bit $iE$.  All other predicate bits are nonsignificant for that operation.
Unless an instruction explicitly produces a complete predicate image or uses
zeroing movement, a false governing bit annuls that lane: it performs no
(:term:base.terms.memory_access:), produces no exception, and preserves the old destination lane.

An instruction that writes a predicate result writes a complete packed
predicate image.  It sets or clears each significant result bit and clears
all nonsignificant bits.  Integer comparisons support \texttt{EQ},
\texttt{NE}, \texttt{ULT}, \texttt{UGE}, \texttt{ULE}, \texttt{UGT},
\texttt{LT}, \texttt{GE}, \texttt{LE}, and \texttt{GT}. VECTORFP defines its
own floating-point comparison forms.
(:ref:VECTOR.instructions.VTESTZ:) and (:ref:VECTOR.instructions.VTESTNZ:) are distinct integer test operations.
Vector integer operations do not update scalar \texttt{FLAGS}.

Integer reductions return a zero-extended result in an \texttt{Rn}.  For an
empty predicate, ADD, OR, and XOR return zero; AND and unsigned MIN return the
all-ones value of the selected width; signed MIN returns the signed maximum;
signed MAX returns the signed minimum; and unsigned MAX returns zero.

(:ref:VECTOR.instructions.PLOOP:) consumes a remaining-count register and an offset register.  If
the remaining count is zero, it takes its encoded branch target and does not
change the predicate or either register.  Otherwise, an offset greater than
or equal to the selected lane count raises (:ref:VECTOR.events.EXCEPTION.VECTOR_LANE_INDEX_OUT_OF_RANGE:) with
error code zero and leaves the instruction uncommitted.  A valid nonzero
iteration activates the next
$\min(remaining, lane\_count-offset)$ lanes, subtracts that number from the
remaining count, clears the offset register, and falls through.  Software
advances the offset between chunks when it needs a nonzero subsequent chunk
origin.
